\begin{abstract}
Although large-scale unsupervised language models (LMs) acquire extensive factual knowledge and exhibit some capability for reasoning, exerting fine-grained control over their outputs remains challenging, primarily due to the inherent lack of supervision in their training processes. To address this limitation, current strategies collect human-annotated rankings of generated outputs and adapt the base LM via fine-tuning to better match these preferences, frequently employing reinforcement learning from human feedback (RLHF). However, RLHF typically involves a complicated and sometimes unstable workflow, requiring first the training of a reward model that encapsulates human preferences, followed by reinforcement learning to optimize the original LM according to this learned reward—without deviating excessively from the pretrained parameters. In this work, we propose a novel reward model parameterization within the RLHF framework, making it possible to obtain the corresponding optimal policy analytically, thereby reformulating the conventional RLHF challenge as an optimization with a straightforward classification loss. The proposed method, termed \textit{Direct Preference Optimization} (DPO), provides a robust, effective, and computationally efficient alternative, obviating the necessity for LM sampling during fine-tuning and greatly reducing the dependence on extensive hyperparameter searches. Empirical results demonstrate that DPO enables LMs to be tuned to human preferences at least as effectively as, and in many cases more effectively than, prior approaches. In particular, DPO outperforms PPO-based RLHF in sentiment steering, and matches or surpasses prior methods in both summarization and single-turn dialogue quality, all while maintaining significant implementation and training simplicity.
\end{abstract}

\section{Introduction}
Large-scale unsupervised language models (LMs), when trained on massive corpora, have demonstrated unexpected and remarkable abilities~\citep{chowdhery2022palm, brown2020language, touvron2023llama,bubeck2023sparks}. Nevertheless, the data underpinning these models originate from human sources encompassing a broad spectrum of objectives, priorities, and competencies. Not all of these human traits are ideal to replicate; for instance, it is advantageous for an AI coding assistant to \textit{recognize} frequent programming errors so as to correct them, but when generating new code, it is preferable for the model to emulate the (possibly infrequent) high-level programming skill present in parts of the dataset. In an analogous manner, a language model should \textit{know about} widespread misconceptions—such as those endorsed by 50\% of individuals—while being prevented from reproducing these inaccuracies as factual in half of the relevant responses. Consequently, it is vital to differentiate and select the model's \emph{targeted responses and behavior} from its vast \textit{repository of knowledge and skills}, to ensure that AI systems are reliable, effective, and under human control \citep{ouyang2022training}. Although prevailing approaches often align LMs with human preferences via reinforcement learning (RL), our findings reveal that the objective underpinning RL-based techniques can be precisely optimized using a straightforward binary cross-entropy loss, thus streamlining the entire preference learning process.

In general terms, current methodologies inject intended behaviors into language models by leveraging carefully constructed datasets containing human preferences that illustrate desirable, safe, and helpful behaviors. This stage of preference learning takes place after the model has undergone large-scale unsupervised pre-training using extensive text corpora. The simplest method to teach preferences is to employ supervised fine-tuning with high-quality human demonstrations. However, the predominant and most effective strategies involve reinforcement learning from human or AI feedback (RLHF/RLAIF; \citep{christiano2017deep,bai2022constitutional}). In RLHF, a reward model is trained using datasets containing human preference judgments, and reinforcement learning is subsequently used to optimize the language model's policy, encouraging it to generate outputs that score highly with respect to the reward, all while constraining deviation from the base model. Although RLHF has led to impressive improvements in conversational fluency and code generation, it introduces significant complexity relative to supervised approaches, as it requires training multiple language models and conducting policy sampling throughout the RL training loop, thus demanding substantial computational resources.

This study presents a method for aligning language models with human preferences without requiring explicit reward models or reinforcement learning techniques. We introduce \textit{{Direct Preference Optimization} (DPO)}, a novel approach that inherently maximizes the same reward objective as traditional RLHF frameworks—which includes a KL-divergence regularization—yet features a more accessible implementation and simplified training process. The core mechanism behind DPO involves increasing the relative log likelihood of preferred responses compared to less desirable ones, but it uniquely incorporates a dynamically adjusted, example-specific importance weight. This weighting addresses the issue of model degradation that arises with a basic probability ratio formulation. Similar to previous techniques, DPO is grounded in a theoretical preference model (for example, the Bradley-Terry model; \cite{bradley1952rankanalysis}), which assesses the congruence between a reward function and observed preference annotations. Distinctly, whereas conventional pipelines employ the preference model to establish a preference loss for training a reward model—followed by optimizing a policy against this learned reward—DPO employs a variable substitution to reformulate the preference loss in terms of the policy parameters directly. As a result, when supplied with a dataset of human-annotated preferences on generated responses, DPO facilitates policy optimization via a straightforward binary cross entropy loss, yielding a policy that is optimal with respect to an implicit reward function fitted to the preference data.

The principal contribution of this work is Direct Preference Optimization (DPO): a robust, reinforcement learning-free strategy for preference-based language model training. Our empirical results demonstrate that DPO matches or surpasses the performance of established methods, including RLHF variants based on PPO, across several tasks involving preference learning, such as sentiment manipulation, text summarization, and conversational modeling, on language models containing up to 6B parameters.

\section{Related Work}

Large-scale self-supervised language models have demonstrated the capacity to tackle some tasks without explicit training examples (zero-shot) \citep{radford2019language} or when provided with a small number of prompts (few-shot) \citep{gpt3,megatron,chowdhery2022palm}. Nonetheless, refining these models with instruction-based datasets containing human-authored completions markedly enhances their task performance and better aligns outputs with user intentions \citep{mishra-etal-2022-cross,sanh2022multitask,chung2022scaling,thoppilan2022lamda}. This approach, known as instruction-tuning, enables large language models (LLMs) to generalize more effectively to novel instructions not present during the tuning phase, thus improving their practicality and versatility \citep{chung2022scaling}. While instruction tuning has achieved impressive outcomes, gathering comparative human evaluations of output quality tends to be more feasible than obtaining expert-crafted examples. Consequently, subsequent research has explored fine-tuning LLMs on collections of human preference data, leading to advancements in areas such as translation \citep{kreutzer-etal-2018-reliability}, summarization \citep{stiennon2022learning,ziegler2020finetuning}, story generation \citep{ziegler2020finetuning}, and following user instructions \citep{ouyang2022training,ramamurthy2023is}. Typically, these strategies begin by training a neural network-based reward model aligned with the collected preference data, employing preference models such as the Bradley-Terry model \citep{bradley1952rankanalysis}. Subsequently, the language model is further tuned to maximize these learned rewards through reinforcement learning methods—most notably, REINFORCE \citep{williams1992reinforce}, proximal policy optimization (PPO; \cite{schulman2017proximal}), or their extensions \citep{ramamurthy2023is}. Parallel advancements have used LLMs, adapted for instruction adherence via human feedback, to create supplementary synthetic preference datasets targeting dimensions like safety or harmlessness \citep{bai2022constitutional}; here, only minimal human supervision—typically a textual rubric guiding LLM annotations—is required. These techniques unite research on applying reinforcement learning to language model training for diverse goals~\citep{Ranzato2015SequenceLT,paulus2018a,wu2018learning} with general strategies for leveraging human preferences in machine learning \citep{christiano2017deep,kupcsik2018learning}. Although utilizing relative human preferences is attractive, practical limitations persist for fine-tuning LLMs with reinforcement learning; the present work introduces a theoretically principled method for preference optimization that circumvents the need for RL.

Beyond linguistic applications, the study of learning policies from preferences has been explored in both contextual bandit and reinforcement learning frameworks, with various methodologies developed. In contextual bandit scenarios, learning that utilizes preference or ranking information instead of reward signals is referred to as the contextual dueling bandit (CDB; \cite{yue2012karmed,dudik2015contextual}). For these settings lacking explicit reward feedback, the theoretical framework replaces the concept of the optimal policy with that of the \textit{von Neumann winner}, defined as a policy whose expected likelihood of outperforming any alternative policy is at least 50\% \citep{dudik2015contextual}. Unlike the CDB paradigm—where feedback on preferences is provided sequentially in an online fashion—human preference learning typically operates on a static dataset of pre-annotated preference-labeled action pairs \citep{yan2022human}. In parallel, \textit{preference-based RL} (PbRL) methodologies train agents based on binary preference feedback supplied by an \textit{unknown} scoring function, as opposed to explicit reward signals \citep{BusaFekete2014,ruiz2023dueling}. There is a variety of PbRL algorithms, some of which can leverage off-policy preference data, but most of these require the prior estimation of the underlying scoring (reward) model before policy optimization proceeds \citep{jain2013learning,BusaFekete2014,christiano2017deep,sadigh2017active,kupcsik2018learning}. Our contribution departs from this paradigm: we introduce a unified policy learning method that directly adjusts the policy to fulfill observed preferences, circumventing separate reward modeling.

\section{Preliminaries}\label{section:prelims}

We outline the RLHF pipeline as described by \citeauthor{ziegler2020finetuning} and further developed in works such as \citep{stiennon2022learning, bai2022training, ouyang2022training}. The standard procedure is composed of three sequential stages: 1) supervised fine-tuning (SFT); 2) preference collection and reward modeling; and 3) reinforcement learning (RL) policy optimization.

\textbf{SFT}: The RLHF process generally starts by fine-tuning a large, pre-trained language model (LM) using supervised learning on a set of high-quality examples related to the target task(s) (such as dialogue generation or summarization), yielding a model $\pi^\text{SFT}$.

\textbf{Reward Modelling Phase}: The second step involves using the SFT-tuned model to respond to prompts $x$ by generating pairs of completions $(y_1, y_2)\sim \pi^\text{SFT}(y \mid x)$. These pairs are then evaluated by human annotators, who provide their preferences, indicated as $y_w\succ y_l \mid x$, with $y_w$ and $y_l$ denoting, respectively, the favored and unfavored answer from the pair $(y_1, y_2)$. These observed preferences are presumed to arise from an unobserved reward mechanism $r^*(y, x)$. Various preference modeling techniques have been utilized, with the Bradley-Terry (BT) model \cite{bradley1952rankanalysis} serving as a widely adopted option (although broader models such as Plackett-Luce \citep{plackett1975analysis, luce2012individual} can also be used when ranked orderings among multiple completions are available). The BT model expresses the human preference distribution $p^*$ as:

Given a static dataset of pairwise preferences $\mathcal{D}=\bigl\{x^{(i)}, y_w^{(i)}, y_l^{(i)}\bigr\}_{i=1}^N$, with samples drawn from $p^*$, we define a parameterized reward function $r_{\phi}(x, y)$ and learn its parameters using maximum likelihood estimation. By viewing the task as binary classification, the associated objective becomes the negative log-likelihood loss:

\begin{equation}\label{eq:reward_model}
    \mathcal{L}_R(r_{\phi}, \mathcal{D}) = -\mathbb{E}_{(x, y_w, y_l)\sim \mathcal{D}}\left[\log \sigma(r_{\phi}(x, y_w) - r_{\phi}(x, y_l))\right]
\end{equation}

with $\sigma$ denoting the logistic function. For language modeling applications, $r_{\phi}(x, y)$ is typically initialized with parameters from the SFT model $\pi^\text{SFT}(y \mid x)$, augmented with a linear readout atop the last transformer layer that outputs a single reward value \cite{ziegler2020finetuning}. In order to reduce the variance in the reward estimates, prior literature normalizes the rewards so that $\mathbb{E}_{x,y\sim \mathcal{D}}\left[r_\phi(x, y)\right]=0$ for every $x$.

\textbf{RL Fine-Tuning Phase}: During reinforcement learning, the trained reward function serves as a feedback signal for the language model. Following methodologies introduced in earlier works~\citep{jaques2017sequence, jaques2020human}, the learning objective is formulated as

\begin{equation}\label{eq:RL}
\max_{\pi_{\theta}}  \mathbb{E}_{x\sim \mathcal{D}, y\sim \pi_{\theta}(y \mid x)}\left[r_{\phi}(x, y)\right] - \beta\mathbb{D}_{\textrm{KL}}\left[\pi_{\theta}(y\mid x) \| \pi_\text{ref}(y\mid x)\right],
\end{equation}

where the parameter $\beta$ regulates the degree of divergence allowed from the reference policy $\pi_\text{ref}$, which typically refers to the SFT model $\pi^\text{SFT}$. The language model policy $\pi_\theta$ is also initialized with weights from $\pi^\text{SFT}$. Incorporating this regularization term is crucial to ensure that the model does not drift far from the data distribution on which the reward function yields reliable signals, thereby mitigating risks such as reward over-optimization or reduced output diversity. Since language generation is discrete and non-differentiable, reinforcement learning techniques are necessary for optimization. Conventionally \citep{ziegler2020finetuning, stiennon2022learning, bai2022training, ouyang2022training}, the composite reward is defined as $r(x, y) = r_{\phi}(x, y) - \beta (\log \pi_{\theta}(y\mid x) - \log \pi_\text{ref}(y\mid x))$, and the policy is updated via PPO \cite{schulman2017proximal}.

\section{Direct Preference Optimization}\label{sec:DPO}

Seeking to circumvent the complexity of applying RL methods for large-scale language model fine-tuning, we propose a streamlined alternative for policy optimization grounded in direct preference data. Contrary to the conventional RLHF pipeline—which involves learning an explicit reward model followed by RL optimization—our methodology exploits a specific form of reward parameterization, enabling the determination of the optimal policy in a closed-form solution, thereby eliminating the RL loop. As detailed in the subsequent section, our main contribution lies in an analytical correspondence between reward functions and their corresponding optimal policies, allowing us to translate an objective formulated over rewards into one over policies directly. This reparameterization enables end-to-end optimization aligned with preference-based models such as the Bradley-Terry model, without requiring a standalone reward network. In this formulation, the policy model inherently captures both the behavior policy and the (implicit) reward structure.

\textbf{Deriving the DPO objective.} To begin, we consider the standard RL objective from earlier work, Eq.~\ref{eq:RL}, with an arbitrary reward function $r$. Consistent with established results~\citep{peters2007reinforcement, peng2019advantage, korbak2022reinforcement, go2023aligning}, it can be easily verified that the optimal policy for the KL-regularized reward maximization problem in Eq.~\ref{eq:RL} has the structure:

\begin{equation}\label{eq:op_policy}
    \pi_r(y\mid x) = \frac{1}{Z(x)}\pi_\text{ref}(y\mid x)\exp\left(\frac{1}{\beta}r(x, y)\right),
\end{equation}

where the normalization constant, $Z(x) =\sum_{y}\pi_\text{ref}(y\mid x)\exp\left(\frac{1}{\beta}r(x, y)\right)$, serves as the partition function. The full details of the derivation can be found in Appendix \ref{app:derivation1}. Even with access to an MLE-based approximation $r_\phi$ of the true reward $r^*$, direct computation of the partition function $Z(x)$ is computationally intensive~\citep{korbak2022reinforcement, go2023aligning}, which hinders the practicality of using this formulation. Nonetheless, Eq.~\ref{eq:op_policy} can be re-expressed to solve for the reward in terms of the associated optimal policy $\pi_r$, reference policy $\pi_\text{ref}$, and the undetermined partition function $Z(\cdot)$. By applying the logarithm to both sides of Eq.~\ref{eq:op_policy} and manipulating the result, we find:

\begin{equation}\label{eq:main_eq}
    r(x,y) =\beta \log \frac{\pi_r(y\mid x)}{\pi_\text{ref}(y\mid x)} + \beta \log Z(x).
\end{equation}

This reparameterization holds for both the ground-truth reward $r^*$ and its corresponding optimal policy $\pi^*$. Conveniently, the Bradley-Terry model’s reliance on reward differences between outcome pairs (i.e., ${p^*(y_1 \succ y_2 \mid x) = \sigma(r^*(x, y_1) - r^*(x, y_2))}$) allows us, when substituting Eq.~\ref{eq:main_eq} into Eq.~\ref{eq:bradley-terry} for $r^*(x,y)$, to eliminate the partition function. Thus, the preference probability based on human data can be reformulated purely as a function of the optimal policy $\pi^*$ and the reference policy $\pi_\text{ref}$. Therefore, the optimal RLHF policy $\pi^*$, as characterized by the Bradley-Terry model, fulfills the following preference relation:

\begin{equation}\label{eq:objective}
    p^*(y_1\succ y_2 \mid x)=\frac{1}{1 + \exp\left(\beta \log \frac{\pi^*(y_2\mid x)}{\pi_\text{ref}(y_2\mid x)} - \beta \log \frac{\pi^*(y_1\mid x)}{\pi_\text{ref}(y_1\mid x)}\right)}
\end{equation}

A more detailed derivation appears in Appendix~\ref{app:derivation2}. While Eq.~\ref{eq:objective} specifically utilizes the Bradley-Terry model, a comparable procedure can be extended to more general models such as Plackett-Luce~\citep{plackett1975analysis, luce2012individual}, as presented in Appendix~\ref{app:plackett_luce_models}.

Having rewritten the probability of human preferences directly in terms of the optimal policy, we are now in a position to define a maximum likelihood estimation objective for any parameterized policy $\pi_\theta$. Mirroring the approach used in reward modeling (see Eq.~\ref{eq:reward_model}), we arrive at the following policy optimization objective:

\begin{equation}\label{eq:optimum_model}
    \mathcal{L}_\text{DPO}(\pi_{\theta}; \pi_\text{ref}) = -\mathbb{E}_{(x, y_w, y_l)\sim \mathcal{D}}\left[\log \sigma \left(\beta \log \frac{\pi_{\theta}(y

In this formulation, we capture an implicit reward through a different parameterization, where the optimal policy directly corresponds to $\pi_\theta$. Additionally, because our approach mirrors the reparametrized Bradley-Terry model, it inherits desirable theoretical guarantees, such as consistency under certain assumptions about the preference data distribution \cite{bong2022generalized}. We provide a more thorough examination of these theoretical aspects of DPO, and how they relate to prior literature, in Section~\ref{sec:theory}.

\textbf{Mechanics of the DPO update.} To clarify the inner workings of DPO, it is instructive to inspect the gradient of the DPO loss, $\mathcal{L}_\text{DPO}$. Specifically, the gradient with respect to model parameters $\theta$ can be expressed as:

\begin{multline*}\label{eq:gradient}
    \nabla_\theta \mathcal{L}_\text{DPO}(\pi_\theta;\pi_\text{ref}) = \\ -\beta\mathbb{E}_{(x, y_w, y_l) \sim \mathcal{D}} \bigg[\underbrace{\sigma(\hat{r}_\theta(x, y_l) - \hat{r}_\theta (x, y_w))}_\text{greater emphasis when reward ranking is inaccurate}\bigg[\underbrace{\nabla_\theta\log \pi(y_w \mid x)}_\text{push up $y_w$'s probability} - \underbrace{\nabla_\theta\log\pi(y_l \mid x)}_\text{push down $y_l$'s probability}\bigg]\bigg],
\end{multline*}

where the implicit reward, $\hat{r}_\theta(x, y) = \beta \log \frac{\pi_\theta(y \mid x)}{\pi_\text{ref}(y \mid x)}$, is determined by comparing the current policy $\pi_\theta$ and a reference policy $\pi_\text{ref}$ (with further elaboration in Section~\ref{sec:theory}). Conceptually, this gradient guides the model to assign greater probability to completions preferred by the dataset ($y_w$) while reducing the probability assigned to less preferred alternatives ($y_l$). Crucially, the update is modulated by how much the implicit reward model overvalues the less preferred completion relative to the more preferred one, as measured and scaled by $\beta$—reflecting the error in ranking and the intensity of the KL regularization. Empirically, we find that this adaptive weighting is essential, as omitting it (i.e., removing the weighting factor) can cause model collapse, as shown in Appendix Table~\ref{tab:unlikelihood_generations}.

\textbf{DPO overview.} The typical DPO workflow includes the following steps: 1) For each prompt $x$, generate response samples $y_1, y_2 \sim \pi_\text{ref}(\cdot \mid x)$, then annotate these with human preference labels to assemble an offline preference dataset $\mathcal{D} = \{x^{(i)}, y_w^{(i)}, y_l^{(i)}\}_{i=1}^N$; 2) Train the language model $\pi_\theta$ to minimize $\mathcal{L}_\text{DPO}$ using the specified $\pi_\text{ref}$, the constructed dataset $\mathcal{D}$, and the intended $\beta$ parameter. In real-world applications, practitioners often prefer to leverage existing public preference datasets rather than create new samples and collect human annotations themselves. Because these datasets are generally generated using $\pi^\text{SFT}$, whenever possible we set $\pi_\text{ref} = \pi^\text{SFT}$. If $\pi^\text{SFT}$ is unavailable, $\pi_\text{ref}$ is instead initialized by maximizing the likelihood of preferred outputs ${(x, y_w)}$, specifically, ${\pi_\text{ref} = \argmax_{\pi}\mathbb{E}_{x, y_w \sim \mathcal{D}}\left[\log \pi(y_w \mid x)\right]}$. This approach serves to reduce the impact of distributional mismatches between the ideal but inaccessible reference distribution and the chosen $\pi_\text{ref}$ employed in DPO. Additional information about the implementation details and hyperparameter choices is provided in Appendix~\ref{app:implementation}.

\section{Theoretical Analysis of DPO} Here, we further interpret the DPO framework, offering theoretical justification for its design, and examine how the strengths of DPO address limitations found in RLHF actor-critic algorithms, such as PPO~\cite{schulman2017proximal}.

\label{sec:theory}

\subsection{Your Language Model Is Secretly a Reward Model} DPO uniquely avoids both explicit reward modeling and reinforcement learning procedures by relying on a unified maximum likelihood approach. The optimization target in Eq. \ref{eq:main_eq} corresponds to a Bradley-Terry formulation, where the reward is parameterized as $r^*(x, y) = \beta \log\frac{\pi^*_\theta(y \mid x)}{\pi_\text{ref}(y \mid x)}$. Our task becomes the optimization of the parametric model $\pi_\theta$, effectively mirroring the reward model objective in Eq. \ref{eq:reward_model} once variables are transformed. This section will develop the theoretical rationale for such a reparameterization, demonstrate that it places no additional constraints on the reward function class, and show that it enables exact recovery of the optimal policy. We start by formalizing an equivalence relation among reward functions.

\begin{definition}
We define two reward functions $r(x, y)$ and $r'(x, y)$ to be equivalent if and only if there exists a function $f$ such that ${r(x, y)-r'(x, y) = f(x)}$.
\end{definition}

It is straightforward to verify that this defines an equivalence relation, which organizes the set of all reward functions into equivalence classes. The following two lemmas hold under this framework:

\begin{lemma}\label{lemma:same_prefrence} In the context of the Plackett-Luce model—and by extension, the Bradley-Terry model—any two reward functions belonging to the same equivalence class induce identical preference distributions.
\end{lemma}

\begin{lemma}\label{lemma:same_policy}
    Any pair of reward functions within the same equivalence class result in the same optimal policy for the associated constrained reinforcement learning problem.
\end{lemma}

We omit the simple proofs here and instead refer the reader to Appendix \ref{app:lemma1}. The first lemma echoes the well-documented issue of under-specification found in the Plackett-Luce family of preference models \cite{plackett1975analysis}. This ambiguity necessitates introducing extra constraints to make the maximum likelihood estimate from Eq.~\ref{eq:reward_model} well-defined \cite{bong2022generalized}. The second lemma asserts that every reward function within an equivalence class gives rise to the same optimal policy, so, when it comes to our learning goal, we only need to recover one representative from the target equivalence class. We establish the following theorem, with the proof presented in Appendix~\ref{app:thm1}:

\begin{theorem}\label{thm:main}
    Provided certain mild assumptions hold, every equivalence class of reward functions consistent with the Plackett-Luce (including Bradley-Terry) model admits a representative of the form ${r(x, y) = \beta \log \frac{\pi(y\mid x)}{\pi_\text{ref}(y\mid x)}}$ for some candidate model $\pi(y\mid x)$ and some reference model $\pi_\text{ref}(y \mid x)$.
\end{theorem}

\begin{sproof}
    Take any reward function $r(x, y)$ and let $\pi_r(y \mid x)$ be its associated optimal model, as characterized in Eq.~\ref{eq:op_policy}. We will demonstrate that there exists a reward function in the equivalence class of $r$ that takes the reparameterized form above. Let us define the projection operator $f$ by
\begin{equation}
    f(r; \pi_\text{ref}, \beta)(x, y) = r(x, y) - \beta\log\sum_{y}\pi_\text{ref}(y\mid x)\exp\left(\frac{1}{\beta}r(x, y)\right)
\end{equation}
Here, $f$ serves to normalize the reward by subtracting the log-partition function for $\pi_r$. Since this normalization term depends only on $x$, $f(r; \pi_\text{ref}, \beta)(x, y)$ remains within the equivalence class of $r(x, y)$. Now, substituting $r$ using the right-hand side of Eq.~\ref{eq:main_eq}, which is valid for any reward function, yields $f(r; \pi_\text{ref}, \beta)(x, y) = \beta \log \frac{\pi_r(y\mid x)}{\pi_\text{ref}(y\mid x)}$. Thus, the operator $f$ constructs a representative of $r$’s equivalence class in the desired reparameterized form, confirming that this approach preserves full generality for modeling rewards.
\end{sproof}

An alternative interpretation of Theorem~\ref{thm:main} is that it identifies the specific reward function, within each equivalence class, that the DPO reparameterization yields. Specifically, this is the reward function that satisfies the following condition:

\begin{equation}\label{eq:lag_p}
     \sum_{y}\underbrace{\pi_\text{ref}(y\mid x)\exp\left(\frac{1}{\beta}r(x, y)\right)}_{=\pi(y\mid x)\text{, using Thm.~\ref{thm:main} reparam.}} = 1,
\end{equation}

which ensures that $\pi(y\mid x)$ forms a proper probability distribution (i.e., all probabilities are non-negative and sum to one). Notably, Eq.~\ref{eq:lag_p} serves as the partition function for the optimal policy that arises from the reward function $r(x, y)$, as seen by comparing to Eq.~\ref{eq:op_policy}. The principal insight behind the DPO approach lies in strategically constraining the otherwise underspecified Plackett-Luce (and, by extension, Bradley-Terry) preference model families. This allows retention of the range of expressible reward models, while explicitly guaranteeing that the optimal policy specified in Eq.~\ref{eq:op_policy} is tractable in closed form for every prompt $x$.

\subsection{Instability of Actor-Critic Algorithms} The proposed framework also facilitates an analysis of instabilities that arise with conventional actor-critic techniques in RLHF, such as PPO. Focusing on the RL fine-tuning stage as presented in Section \ref{section:prelims}, parallels can be drawn with the control as inference paradigm \cite{levine2018reinforcement} when addressing the constrained RL formulation stated in \ref{eq:RL}. Here, we parameterize the model as $\pi_{\theta}(y\mid x)$ and seek to minimize the KL divergence $\mathbb{D}_{\text{KL}}[\pi_{\theta}(y|x) \mid\mid \pi^*(y\mid x)]$, where $\pi^*$ corresponds to the optimal policy defined by Eq.~\ref{eq:optimum_model} in terms of a learned reward $r_{\phi}(y, x)$. With straightforward manipulation, this yields the following optimization problem:

\begin{equation}\label{eq:AC}
    \max_{\pi_{\theta}}\mathbb{E}_{\pi_{\theta}(y\mid x)}\left[\underbrace{r_{\phi}(x, y) -\beta\log\sum_{y}\pi_\text{ref}(y\mid x)\exp\left(\frac{1}{\beta}r_{\phi}(x, y)\right)}_{f(r_{\phi}, \pi_\text{ref}, \beta)} - \underbrace{\beta\log\frac{\pi_{\theta}(y\mid x)}{\pi_\text{ref}(y\mid x)}}_{\text{KL}}\right]
\end{equation}

The objective optimized here aligns with that of previous research 
\citep{ziegler2020finetuning, stiennon2022learning, bai2022training, ouyang2022training}, employing a DPO-equivalent reward within the reward class $r_{\phi}$. In this framework, the normalization component in $f(r_{\phi}, \pi_\text{ref}, \beta)$ can be viewed as the soft value function for the reference policy $\pi_\text{ref}$. Although this normalization does not influence the final optimal policy, omitting it may result in high-variance policy gradients and thus hamper the stability of learning. This normalization can be managed by leveraging a learned value function, though optimizing it is itself a nontrivial task. An alternative strategy observed in earlier studies involves normalizing rewards against a human completion baseline, which effectively provides a single-sample Monte-Carlo approximation of the normalization term. Notably, the DPO reparameterization produces a reward function that does not rely on any baselines.

\section{Experiments}
Here, we present an empirical assessment of DPO’s capability to train policies based solely on preference data. We begin with a controlled text-generation scenario, evaluating how DPO balances reward maximization and KL-divergence minimization with respect to the reference policy when compared to established preference optimization methods such as PPO. Subsequently, we analyze DPO’s effectiveness with larger-scale models and on more challenging RLHF applications, including both summarization and dialogue tasks. Our experiments show that, even with minimal hyperparameter tuning, DPO generally matches or exceeds the performance of robust baselines like RLHF with PPO, and outperforms the approach of selecting the best out of $N$ sampled trajectories under a trained reward model. The following sections outline our experimental methodology; further experimental specifics are available in Appendix~\ref{app:exp_details}.

\textbf{Tasks.} Our study investigates three distinct open-ended text generation tasks. For each, policies are trained using a preference dataset $\mathcal{D}=\bigl\{x^{(i)}, y_w^{(i)}, y_l^{(i)}\bigr\}_{i=1}^N$. In the \textbf{controlled sentiment generation} setting, $x$ corresponds to the opening of a movie review sampled from the IMDb dataset \cite{maas-EtAl:2011:ACL-HLT2011}, and the objective is for the policy to produce a continuation $y$ exhibiting positive sentiment. To facilitate a controlled assessment, we automatically generate preference pairs from model outputs using a pre-trained sentiment classifier, where $p(\text{positive}\mid x,y_w)>p(\text{positive}\mid x,y_l)$. For SFT, we train GPT-2-large on the IMDb training set until performance plateaus (additional information provided in App~\ref{app:sentiment_details}). In the \textbf{summarization} task, $x$ is a Reddit forum post, and the policy must write a concise summary $y$ capturing the core content. We utilize the Reddit TL;DR dataset \citep{volske-etal-2017-tl} as well as human preferences obtained by \citeauthor{stiennon2022learning}, following conventions in prior studies. Our SFT model is trained on summaries written by humans\footnote{\url{https://huggingface.co/CarperAI/openai_summarize_tldr_sft}} and optimized with the TRLX \citep{leandro_von_werra_2023_7790115} library for RLHF. The preference labels derive from \citeauthor{stiennon2022learning}’s evaluations over outputs from a similarly fine-tuned, yet distinct, SFT model. The third task, \textbf{single-turn dialogue}, involves $x$ as a single user query—ranging from factual inquiries to personal advice requests. The aim is for the policy to return an informative and engaging answer $y$. For this, we use the Anthropic Helpful and Harmless dataset \citep{bai2022training}, consisting of 170,000 dialogues between a user and an automated system. Each dialogue concludes with a pair of responses produced by a large language model, and a label denoting which response a human prefers. Since a pre-trained SFT model is not available in this scenario, we construct the SFT model by fine-tuning a general-purpose language model using only responses corresponding to the preferred completions.

\textbf{Evaluation.} We adopt two complementary evaluation strategies in our experiments. To investigate each algorithm’s capability in optimizing the constrained reward maximization objective within the controlled sentiment generation setting, we characterize performance by plotting the trade-off between obtained reward and KL-divergence relative to the reference policy. Since access to the exact reward function (a sentiment classifier) allows calculation of this frontier, this metric provides clear comparisons. In contrast, real-world scenarios lack a known ground-truth reward function. Accordingly, for such cases, we utilize a \textit{win rate} metric against a baseline policy, leveraging GPT-4 as a surrogate for human judgments. Specifically, this proxy assesses summary quality and helpfulness of dialogue responses in the summarization and single-turn dialogue tasks. Baselines are defined as reference summaries from the test set for summarization and as the preferred responses from the test dataset for dialogue. While research indicates language models can surpass automated metrics for evaluation \citep{Chen2023ExploringTU}, we corroborate the validity of GPT-4-based evaluation through a targeted human study described in Sec.~\ref{sec:human-judgments}. Our findings reveal strong alignment between GPT-4 and human judgments, with the level of human-GPT-4 agreement frequently matching or exceeding the agreement between different human annotators.

\textbf{Methods.} Alongside DPO, we assess multiple established techniques for training language models to better reflect human preferences. As straightforward baselines, we apply zero-shot prompting of \textbf{GPT-J} \citep{gpt-j} in the summarization setting, and 2-shot prompting with \textbf{Pythia-2.8B} \citep{biderman2023pythia} for the dialogue task. Beyond prompting, we evaluate the \textbf{SFT} model, and the \textbf{Preferred-FT} model, which results from supervised fine-tuning on the preferred output $y_w$ selected from the SFT model (for controlled sentiment and summarization), or from a general language model (in the case of single-turn dialogue). We further consider the \textbf{Unlikelihood}~\citep{welleck2019neural} strategy, a pseudo-supervised method that promotes higher probability for $y_w$ while penalizing the probability for $y_l$, incorporating an optional coefficient $\alpha\in[0,1]$ to weight the unlikelihood penalty. We also investigate the \textbf{PPO} approach \citep{schulman2017proximal}, utilizing a reward model trained on preference annotations, as well as \textbf{PPO-GT}, which serves as an oracle algorithm trained directly with the ground truth reward in the controlled sentiment context. For sentiment experiments, two PPO-GT implementations are included: a standard implementation \cite{leandro_von_werra_2023_7790115} and an enhanced version that normalizes rewards and adjusts hyperparameters to boost performance (similar adjustments are applied to PPO trained on learned rewards). Lastly, we introduce the \textbf{Best of $N$} approach: $N$ candidate responses are sampled from the SFT model (or Preferred-FT for dialogue), with the highest-scoring completion—evaluated via the reward model learned from preferences—returned as output. Although this method generally achieves strong performance by separating the reward model’s effectiveness from PPO’s learning dynamics, it becomes computationally infeasible for even modest $N$, as every test-time query requires generation of $N$ separate completions.

\subsection{How well can DPO optimize the RLHF objective?}

RLHF methods commonly employ a KL-constrained reward maximization formulation, which ensures that the policy optimizes for higher reward while penalizing large deviations from the reference policy. Thus, comparative evaluations of algorithms should jointly consider both the level of reward attained and the extent of KL divergence; achieving marginally higher reward at the expense of a substantial increase in KL is not ideal. Figure~\ref{fig:frontier-tldr-main} presents the reward-KL tradeoff curves for several algorithms within the sentiment classification task. For each algorithm, we perform several training runs, each varying a hyperparameter that controls the conservativeness of the policy (target KL $\in\{3,6,9,12\}$ for PPO, $\beta \in \{0.05,0.1,1,5\}$ for DPO, $\alpha\in\{0.05,0.1,0.5,1\}$ for unlikelihood, and different random seeds for preferred-FT), leading to a total of 22 distinct experiments. At intervals of 100 training steps up to convergence, each resulting policy is evaluated on a suite of test prompts, where we calculate both the mean reward using the true reward function and the average sequence-level KL\footnote{Specifically, we sum KL divergences at each time step.} with respect to the reference policy, $\text{KL}\left(\pi\mid \mid \pi_\text{ref}\right)$. Our observations indicate that DPO consistently yields a highly favorable reward-KL frontier, surpassing other approaches in maximizing reward while simultaneously maintaining low KL. This outcome is especially significant for a few key reasons: DPO and PPO are trained on the same objective, yet DPO displays a marked improvement in efficiency, dominating the reward/KL tradeoff compared to PPO. Moreover, DPO outperforms PPO even in scenarios where PPO is granted access to the actual ground truth reward signal (PPO-GT).

\subsection{Can DPO scale to real preference datasets?}
\label{sec:dpo-real-datasets}
We proceed by assessing how DPO performs when fine-tuned on real-world datasets for summarization and single-turn dialogue. In the context of summarization, traditional automated metrics like ROUGE often exhibit weak correlation with human preferences~\citep{stiennon2022learning}. Previous research has demonstrated that optimizing language models with PPO based on human preferences results in more satisfactory summaries. To compare various approaches, we generate completions on the TL;DR summarization dataset's test split and calculate the mean win rate relative to reference completions. For all considered methods, outputs are sampled using temperature settings ranging from 0.0 to 1.0, with the corresponding win rates depicted in Figure~\ref{fig:frontier-tldr-main} (right). The DPO, PPO, and Preferred-FT strategies all fine-tune the same GPT-J SFT model\footnote{\url{https://huggingface.co/CarperAI/openai_summarize_tldr_sft}}. Our results indicate that DPO achieves a win rate of about 61\% when sampling at temperature 0.0, outperforming PPO, which attains roughly 57\% at its best temperature setting of 0.0. Additionally, DPO surpasses the $N$ baseline in terms of the highest win rate observed. It should be noted that DPO’s $\beta$ hyperparameter was not thoroughly optimized in our experiments, suggesting the potential for even stronger results. Furthermore, DPO displays greater resilience to changes in sampling temperature compared to PPO, whose performance may decline to that of the original GPT-J model at higher temperatures. Meanwhile, Preferred-FT shows minimal improvement over the SFT model. Section~\ref{sec:human-judgments} presents a direct comparison of DPO and PPO through human evaluations, revealing that DPO completions sampled at temperature 0.25 were favored in 58\% of cases over PPO samples at temperature 0.

For single-turn dialogue evaluation, we assess various approaches using the segment of the test set from the Anthropic HH dataset \citep{bai2022training} that features a single exchange between a human and an assistant. The win rates for each approach are calculated with GPT-4 by comparing them to the preferred completions from the test set, which serve as the ground truth. Since there is no established SFT model tailored for this task, our methodology begins with the pre-trained Pythia-2.8B model. We employ Preferred-FT to fine-tune a reference model on selected completions, ensuring that model outputs remain within its natural distribution, and subsequently apply DPO training. Our analysis also benchmarks against the top-scoring completion among 128 Preferred-FT generations (since performance for the Best of $N$ method stabilizes at $N=128$ on this dataset; refer to Appendix Figure~\ref{fig:best-of-n}), as well as a 2-shot prompt-based version of the Pythia-2.8B model. Notably, DPO either matches or exceeds the best outcomes from other approaches when temperature is optimized. We further examine an RLHF model using PPO, trained on the Anthropic HH dataset \footnote{\url{https://huggingface.co/reciprocate/ppo_hh_pythia-6B}}, provided by a recognized organization \footnote{\url{https://github.com/CarperAI/trlx/tree/main/examples/hh}}, but fail to identify a prompt or sampling temperature that outperforms the original Pythia-2.8B model. Drawing from our TL;DR findings and the observation that both methodologies target the same reward objective, we treat Best of 128 as an approximate surrogate for PPO results. In summary, DPO stands out as the only method that both improves upon the reference completions in the Anthropic HH dataset and remains computationally efficient, rivaling or exceeding the results of the resource-intensive Best of 128 baseline. Figure~\ref{fig:dialogue-main} further demonstrates that DPO reaches optimal performance rapidly.

\subsection{Generalization to a new input distribution}

To better assess the robustness of PPO and DPO under distributional changes, we test both policies—trained in our Reddit TL;DR summarization setup—on out-of-distribution data drawn from the test portion of the CNN/DailyMail corpus \citep{nallapati-etal-2016-abstractive}. We utilize the top-performing sampling temperatures from the TL;DR experiments (0 and 0.25). Results are shown in Table~\ref{tab:ood}. We measured the rate at which GPT-4 selected the model-generated summary over the ground-truth summaries, employing the identical GPT-4 (C) evaluation prompt used previously, with the substitution of ``news article'' for ``forum post.'' On this novel distribution, the DPO policy again achieves a considerably higher performance than PPO. These findings suggest that DPO retains a comparable generalization capacity to PPO, even without leveraging the extra unlabeled Reddit TL;DR inputs that benefit PPO.

\subsection{Validating GPT-4 judgments with human judgments}
\label{sec:human-judgments}
We designed a human evaluation to examine the alignment between GPT-4 assessments and those made by humans, drawing on results from the TL;DR summarization study and employing two distinct GPT-4 prompts. The \textbf{GPT-4 (S)} prompt asks only which summary better encapsulates the core content of the post. The \textbf{GPT-4 (C)} prompt requests the same, but explicitly asks for a judgment on conciseness; we assess this version because with the \textbf{GPT-4 (S)} prompt, GPT-4 demonstrates a preference for summaries that are longer and more redundant compared to human annotators. Full prompt details are provided in Appendix~\ref{app:prompts}. We conduct three sets of pairwise evaluations, featuring the highest (DPO, temp. 0.25), lowest (PPO, temp. 1.0), and an intermediate (SFT, temp. 0.25) performing model, all compared against the PPO baseline with greedy decoding (the optimal PPO temperature). These selections span a range of output qualities; for DPO and PPO-1 we collect several human assessments to allow for agreement analysis, but limited human resources prevent us from doing so for all models. The results show that, with both prompts, GPT-4's preferences align with human choices at rates comparable to human inter-rater agreement, implying that GPT-4 is a viable substitute for human judgment in this setting. We note that win rates obtained using the \textbf{GPT-4 (C)} prompt tend to mirror human judgment more closely, which is why it is adopted as our main evaluation metric in Section~\ref{sec:dpo-real-datasets}. Additional methodological details, including the design of the web interface for raters and a roster of the human evaluators, are available in Appendix~\ref{app:human-study}.

\section{Discussion}
Preference-based learning offers a robust and extensible approach for developing language models that are both effective and aligned with user values. In this work, we have proposed DPO—a straightforward method for preference-driven language model training that does not rely on reinforcement learning. Unlike traditional approaches that reformulate preference learning within the RL paradigm to utilize standard RL methods, DPO instead establishes a direct correspondence between language model policies and reward functions. This allows for optimizing a language model to align with human preferences via a basic cross-entropy loss, circumventing the need for reinforcement learning while preserving expressiveness. Remarkably, DPO achieves performance on par with or exceeding that of established RLHF techniques such as PPO, and does so with minimal hyperparameter adjustment, thereby significantly simplifying the process of training language models with human preferences.

\textbf{Limitations \& Future Work.} Our findings point to several directions deserving further investigation. One open question is the out-of-distribution generalization capability of policies trained with DPO, especially in comparison to those trained via explicit reward functions; preliminary evidence indicates DPO’s generalization is comparable to that of PPO-based approaches, though more exhaustive analysis is necessary. For instance, could DPO-based policies leverage self-labeled data for learning from unlabeled prompts as effectively? Another key avenue concerns the characterization of reward over-optimization within the DPO context, and whether the modest performance dip observed in Figure~\ref{fig:dialogue-main}-right is a reflection of this phenomenon. Additionally, our evaluations were limited to models up to 6B parameters, and scaling DPO to much larger, state-of-the-art models remains an intriguing prospect. In the area of evaluation, our experiments suggest that GPT-4-based win rate assessments are prompt-dependent; thus, optimizing automated evaluation methods represents a promising research area. Beyond the domain of language models, DPO may prove valuable for training generative models in various other modalities.